\documentclass[fontsize=15pt]{jlreq}

%プリアンブル
\usepackage{amsmath}
\usepackage{mathtools}
\pagestyle{empty}
\usepackage[top=15mm, bottom=20mm, left=20mm, right=20mm]{geometry}

\newcommand{\m}[1]{\raisebox{0.1em}{\textcircled{\scriptsize #1}}}



%本文
\begin{document}
\begin{center}
{\LARGE {振り返り}} \\
-天体\m{1}- %単元ごとに変えるか消す
\end{center}
\vspace{1em}
\begin{flushright}
名前：\underline{\hspace{5cm}}
\end{flushright}

% 問題部分
\begin{enumerate}
    \item 地球から見た太陽が，見かけ上約一日で地球の周りを一周することを何というか．\\
    (\hspace{3cm})
    \vspace{2em} 

    \item 太陽のつくりについて，
    \textbf{表面}，\textbf{中心部}，\textbf{黒点}
    をその温度が高い順に並び替えろ．\\
    (\hspace{3cm}$\to$\hspace{3cm}$\to$\hspace{3cm})
    \vspace{2em}

    \item 天体が天頂より南側で子午線を通過することを
    (\hspace{3cm})といい，南中する時刻を南中時刻，
    南中するときの高度を(\hspace{3cm})という．
    \vspace{2em}

    \item 月はどの天体の周りをまわっているか．\\
    (\hspace{3cm})

    \vspace{2em}

    \item 望遠鏡に太陽投影板をセットし，
    太陽の像を投影した．太陽の黒点の位置を時間ごとに記録して
    いったところ，その黒点は時間が経つにつれて東から西へと移動し，
    黒点の形は太陽の像の中心では円形，周辺部(端側)ではだ円形に見えた．
    このことからわかる太陽の特徴を二つ挙げよ．\\
    (\hspace{15cm})
\end{enumerate}

\end{document}